\documentclass[10pt,a4paper]{amsart}
\usepackage{amsmath,amssymb,amsthm,graphicx}
\usepackage{verbatim}
\newenvironment{reference}{\comment}{\endcomment}
%\newenvironment{reference}{}{}
\newenvironment{slogan}{\comment}{\endcomment}
\newenvironment{history}{\comment}{\endcomment}

% For commutative diagrams we use Xy-pic
\usepackage[all]{xy}

% We use 2cell for 2-commutative diagrams.
\xyoption{2cell}
% Theorem environments.
%
\theoremstyle{plain}
\newtheorem{theorem}[subsection]{Theorem}
\newtheorem{proposition}[subsection]{Proposition}
\newtheorem{lemma}[subsection]{Lemma}

\theoremstyle{definition}
\newtheorem{definition}[subsection]{Definition}
\newtheorem{example}[subsection]{Example}
\newtheorem{exercise}[subsection]{Exercise}
\newtheorem{situation}[subsection]{Situation}

\theoremstyle{remark}
\newtheorem{remark}[subsection]{Remark}
\newtheorem{remarks}[subsection]{Remarks}

\numberwithin{equation}{subsection}

% Macros
%
\def\lim{\mathop{\mathrm{lim}}\nolimits}
\def\colim{\mathop{\mathrm{colim}}\nolimits}
\def\Spec{\mathop{\mathrm{Spec}}}
\def\Hom{\mathop{\mathrm{Hom}}\nolimits}
\def\Ext{\mathop{\mathrm{Ext}}\nolimits}
\def\SheafHom{\mathop{\mathcal{H}\!\mathit{om}}\nolimits}
\def\SheafExt{\mathop{\mathcal{E}\!\mathit{xt}}\nolimits}
\def\Sch{\mathit{Sch}}
\def\Mor{\mathop{\mathrm{Mor}}\nolimits}
\def\Ob{\mathop{\mathrm{Ob}}\nolimits}
\def\Sh{\mathop{\mathit{Sh}}\nolimits}
\def\NL{\mathop{N\!L}\nolimits}
\def\CH{\mathop{\mathrm{CH}}\nolimits}
\def\proetale{{pro\text{-}\acute{e}tale}}
\def\etale{{\acute{e}tale}}
\def\QCoh{\mathit{QCoh}}
\def\Ker{\mathop{\mathrm{Ker}}}
\def\Im{\mathop{\mathrm{Im}}}
\def\Coker{\mathop{\mathrm{Coker}}}
\def\Coim{\mathop{\mathrm{Coim}}}

% Boxtimes
%
\DeclareMathSymbol{\boxtimes}{\mathbin}{AMSa}{"02}

%
% Macros for moduli stacks/spaces
%
\def\QCohstack{\mathcal{QC}\!\mathit{oh}}
\def\Cohstack{\mathcal{C}\!\mathit{oh}}
\def\Spacesstack{\mathcal{S}\!\mathit{paces}}
\def\Quotfunctor{\mathrm{Quot}}
\def\Hilbfunctor{\mathrm{Hilb}}
\def\Curvesstack{\mathcal{C}\!\mathit{urves}}
\def\Polarizedstack{\mathcal{P}\!\mathit{olarized}}
\def\Complexesstack{\mathcal{C}\!\mathit{omplexes}}
% \Pic is the operator that assigns to X its picard group, usage \Pic(X)
% \Picardstack_{X/B} denotes the Picard stack of X over B
% \Picardfunctor_{X/B} denotes the Picard functor of X over B
\def\Pic{\mathop{\mathrm{Pic}}\nolimits}
\def\Picardstack{\mathcal{P}\!\mathit{ic}}
\def\Picardfunctor{\mathrm{Pic}}
\def\Deformationcategory{\mathcal{D}\!\mathit{ef}}
\usepackage{cleveref}
\crefname{originallemma}{original lemma}{original lemmas}
\Crefname{originallemma}{Original lemma}{Original lemmas}
\makeatletter
\newlabel{algebra-lemma-exists-dvr}{{00PH}{source}{}{sourcepointer-1}{}}
\newlabel{algebra-lemma-exists-dvr@cref}{{[originallemma][0][]00PH}{[1][0][]source}{}{}{}}
\makeatother
\crefname{originallemma}{original lemma}{original lemmas}
\Crefname{originallemma}{Original lemma}{Original lemmas}
\makeatletter
\newlabel{formal-spaces-lemma-closed-immersion-into-countably-indexed}{{0AQI}{source}{}{sourcepointer-2}{}}
\newlabel{formal-spaces-lemma-closed-immersion-into-countably-indexed@cref}{{[originallemma][0][]0AQI}{[1][0][]source}{}{}{}}
\makeatother
\crefname{originalreference}{original reference}{original references}
\Crefname{originalreference}{Original reference}{Original references}
\makeatletter
\newlabel{lemma-base-change-rig-surjective}{{0AQS}{source}{}{sourcepointer-3}{}}
\newlabel{lemma-base-change-rig-surjective@cref}{{[originalreference][0][]0AQS}{[1][0][]source}{}{}{}}
\makeatother
\crefname{originalreference}{original reference}{original references}
\Crefname{originalreference}{Original reference}{Original references}
\makeatletter
\newlabel{lemma-monomorphism-rig-surjective}{{0AR0}{source}{}{sourcepointer-4}{}}
\newlabel{lemma-monomorphism-rig-surjective@cref}{{[originalreference][0][]0AR0}{[1][0][]source}{}{}{}}
\makeatother
\crefname{originalreference}{original reference}{original references}
\Crefname{originalreference}{Original reference}{Original references}
\makeatletter
\newlabel{lemma-upshot}{{0AQV}{source}{}{sourcepointer-5}{}}
\newlabel{lemma-upshot@cref}{{[originalreference][0][]0AQV}{[1][0][]source}{}{}{}}
\makeatother
\title{Rig-surjectivity of a closed immersion of formal algebraic spaces}
\author{The Stacks Project Authors}
\date{stacks-project:restricted.tex:local-snapshot-20261012; GNU FDL 1.2 or later}
\begin{document}
\maketitle
\noindent\textit{Editorial note.} All delivered mathematical text is the author's own, copied from the pinned Stacks Project chapter file restricted.tex; no statement or proof step was rewritten, reordered or removed.\par\medskip
\noindent\textit{Editorial note.} The counted proof is the author's complete argument: the proof is complete in both directions\par\medskip
\noindent\textit{Editorial note.} The label lemma-closed-immersion-rig-surjective names the counted statement, and the proof environment immediately following it in the source is delivered verbatim as its proof.\par\medskip
\noindent\textit{Editorial note.} The counted statement is carried by the author's own lemma environment in the chapter source.\par\medskip
\noindent\textit{Allowed standard prerequisites (editorial).}\begin{itemize}
\item closed immersions of locally Noetherian formal algebraic spaces and their taut surjections B to A in WAdm\textasciicircum{}Noeth on affine charts
\item rig-surjectivity, its base change stability, the monomorphism reformulation and the upshot reduction
\item ideals of definition of adic rings and the Jacobson/DVR embedding lemma for one-dimensional quotients
\end{itemize}\medskip
\ifdefined\hypertarget\hypertarget{sourcepointer-1}{}\fi
\noindent\textit{Source reference.} Original lemma 00PH: sources/stacks-chow/tags.txt, source lines 931--931. Prerequisite pointer invoked at proof line(s) 5580; it names a standard fact of a sibling chapter that the counted argument uses, and every step of the counted proof built on it is retained verbatim in the delivered body.\par\medskip
\ifdefined\hypertarget\hypertarget{sourcepointer-2}{}\fi
\noindent\textit{Source reference.} Original lemma 0AQI: sources/stacks-chow/tags.txt, source lines 13511--13511. Prerequisite pointer invoked at proof line(s) 5566; it names a standard fact of a sibling chapter that the counted argument uses, and every step of the counted proof built on it is retained verbatim in the delivered body.\par\medskip
\ifdefined\hypertarget\hypertarget{sourcepointer-3}{}\fi
\noindent\textit{Source reference.} Original reference 0AQS: sources/stacks-restricted/pinned-original-source.tex, source lines 5324--5324. Prerequisite pointer invoked at proof line(s) 5569; it names a standard fact of the same chapter that the counted argument uses, and every step of the counted proof built on it is retained verbatim in the delivered body.\par\medskip
\ifdefined\hypertarget\hypertarget{sourcepointer-4}{}\fi
\noindent\textit{Source reference.} Original reference 0AR0: sources/stacks-restricted/pinned-original-source.tex, source lines 5525--5525. Prerequisite pointer invoked at proof line(s) 5584; it names a standard fact of the same chapter that the counted argument uses, and every step of the counted proof built on it is retained verbatim in the delivered body.\par\medskip
\ifdefined\hypertarget\hypertarget{sourcepointer-5}{}\fi
\noindent\textit{Source reference.} Original reference 0AQV: sources/stacks-restricted/pinned-original-source.tex, source lines 5394--5394. Prerequisite pointer invoked at proof line(s) 5590; it names a standard fact of the same chapter that the counted argument uses, and every step of the counted proof built on it is retained verbatim in the delivered body.\par\medskip
\begin{lemma}
\label{lemma-closed-immersion-rig-surjective}
Let $S$ be a scheme. Let $f : X \to Y$ be a morphism of formal algebraic
spaces. Assume that $X$ and $Y$ are locally Noetherian and $f$ is a
closed immersion. The following are equivalent
\begin{enumerate}
\item $f$ is rig-surjective, and
\item for every affine formal algebraic space $V$ and every morphism
$V \to Y$ which is representable by algebraic spaces and \'etale
the morphism $X \times_Y V \to V$ corresponds to a surjective morphism
$B \to A$ in $\textit{WAdm}^{Noeth}$ whose kernel $J$ has the following
property: $IJ^n = 0$ for some ideal of definition $I$ of $B$
and some $n \geq 1$.
\end{enumerate}
\end{lemma}

\begin{proof}
Let us observe that given $V$ and $V \to Y$ as in (2) without any
further assumption on $f$ we see that the morphism $X \times_Y V \to V$
corresponds to a surjective morphism $B \to A$ in $\textit{WAdm}^{Noeth}$
by Formal Spaces, Lemma
\ref{formal-spaces-lemma-closed-immersion-into-countably-indexed}.

\medskip\noindent
Assume (1). By Lemma \ref{lemma-base-change-rig-surjective} we see that
$\text{Spf}(A) \to \text{Spf}(B)$ is rig-surjective.
Let $I \subset B$ be an ideal of definition. Since $B$ is adic,
$I^m \subset B$ is an ideal of definition for all $m \geq 1$.
If $I^m J^n \not = 0$ for all $n, m \geq 1$, then
$IJ$ is not nilpotent, hence $V(IJ) \not = \Spec(B)$.
Thus we can find a prime ideal $\mathfrak p \subset B$
with $\mathfrak p \not \in V(I) \cup V(J)$.
Observe that $I(B/\mathfrak p) \not = B/\mathfrak p$
hence we can find a maximal ideal
$\mathfrak p + I \subset \mathfrak m \subset B$.
By Algebra, Lemma \ref{algebra-lemma-exists-dvr}
we can find a discrete valuation ring $R$
and an injective local ring homomorphism $(B/\mathfrak p)_\mathfrak m \to R$.
Clearly, the ring map $B \to R$ cannot factor through $A = B/J$.
According to Lemma \ref{lemma-monomorphism-rig-surjective}
this contradicts the fact that $\text{Spf}(A) \to \text{Spf}(B)$
is rig-surjective. Hence for some $n, m$ we do have
$I^n J^m = 0$ which shows that (2) holds.

\medskip\noindent
Assume (2). By Lemma \ref{lemma-upshot} it suffices to show
that $\text{Spf}(A) \to \text{Spf}(B)$ is rig-surjective.
Pick an ideal of definition $I \subset B$ and an integer $n$
such that $I J^n = 0$.
Consider a ring map $B \to R$ where $R$ is a discrete valuation
ring and the image of $I$ is nonzero. Since $R$ is a domain, we
conclude the image of $J$ in $R$ is zero. Hence $B \to R$
factors through the surjection $B \to A$ and we are done by
definition of rig-surjective morphisms.
\end{proof}

\end{document}
